\chapter[Predigt zur Wiederbeziehung des Concordia-College]{Predigt, gehalten bei der Wiederbeziehung des renovirten Concordia-College zu Fort Wayne, Ind., am 11. October 1903}
\quellseite{2}

\begin{center}
(Auf Beschluß des Lehrercollegiums dem Druck überlassen von P. Aug. C. Senne.)
\end{center}

Text: Luc. 24, 29.: „Und sie nöthigten ihn und sprachen: Bleibe bei uns, denn es will Abend werden, und der Tag hat sich geneiget. Und er ging hinein, bei ihnen zu bleiben.“

In Christo, unserm hochgelobten HErrn und Heilande, geliebte Glaubens- und Festgenossen!

Unsere theure Concordia prangt heute im Festschmuck. Und wir, die wir von nah und fern uns hier versammelt haben, sind ebenfalls festlich und fröhlich gestimmt. Ja, in lauten Lob- und Jubelliedern haben wir bereits der Freude und dem Dank unseres Herzens Ausdruck gegeben.

Sagt, was hat das alles zu bedeuten? Worin besteht das festliche Ereigniß, das wir heute mit Lob und Dank gegen Gott begehen? So es betrifft ja unsere hiesige Lehranstalt, unsere theure Concordia, die uns allen so sehr am Herzen liegt. Ein Hauptgebäude derselben, das nach langjährigem Gebrauch an allen Ecken und Enden abgenutzt und morsch geworden war, steht nun wieder gründlich ausgebessert da. Vom Keller bis zum Dach ist es in guten Zustand versetzt worden. Mit Lust und Freude können nun unsere lieben Professoren in den neuhergestellten Räumen wieder unterrichten. Mit Lust und Freude können nun unsere lieben Knaben und Jünglinge in denselben wieder studiren.

Wie, ist das für uns Christen nicht Veranlassung genug zum freudigen Loben und Danken Gottes? Ganz gewiß! So wollen wir es denn auch heute mit Herz und Mund bekennen und rühmen: „Der Herr hat Großes an uns gethan, deß sind wir fröhlich.“ Ja: „Lobe den Herrn, meine Seele, und was in mir ist, seinen heiligen Namen. Lobe den Herrn, meine Seele, und vergiß nicht, was er dir Gutes gethan hat.“

Doch, meine Lieben, wir sollen und wollen heute nicht nur danken. Wir sollen und wollen heute auch beten. Ja, in diesen letzten Zeiten bedarf gerade diese unsere Lehranstalt, in
\quellseite{3}
der unsere Jünglinge ihre Vorbildung für das heilige Predigtamt erhalten, unsers Gebets und des Gebets unserer ganzen Synode gar sehr. Aber was sollen und wollen wir beten? Gewiß, in dem Gebet der Emmausjünger: „Herr, bleibe bei uns!“ haben wir das rechte Gebet für unsere Concordia. Das sei denn auch jetzt der Gegenstand unserer weiteren Andacht, nämlich:

\section{Unser Gebet: Herr Jesu, bleibe bei uns!}

Wir betrachten:

1. wie getrost und zuversichtlich wir dieses Gebet thun können;

2. welch hohe Ursache wir zu diesem Gebet haben.

1.

„Bleibe bei uns!“ so baten jene beiden Jünger den Herrn Jesum, als sie spät am Osternachmittag mit einander in Emmaus angekommen waren. Und wie, konnten sie diese Bitte nicht getrost und zuversichtlich thun? Jesus hatte sich ja so freundlich zu ihnen gestellt, hatte sie begleitet und war bis jetzt bei ihnen geblieben. Dazu hatte er sich ihrer unterwegs so treulich angenommen. Er hatte sie belehrt, ihnen zurechtgeholfen, sie getröstet, sie gestärkt und sie so selig gemacht. O gewiß, was die Emmausjünger von Jesu erfahren hatten, machte sie getrost und zuversichtlich zu bitten: „Bleibe bei uns!“

„Herr Jesu, bleibe bei uns!“ Diese Bitte können auch wir ebenso getrost und zuversichtlich thun wie einst die Emmausjünger, wenn wir auf unsere theure Concordia sehen. Alles, was diese unsere Anstalt bisher von Jesu erfahren hat, berechtigt und ermuntert uns zu dieser Bitte. Jesus ist von Anfang an bis auf den heutigen Tag mit dieser Anstalt gewesen. Durch sein Wort, welches hier wohnte und im Schwange ging, hat er selbst hier gewohnt, gewaltet und gewirkt. Er hat hier seine gnadenreiche Gegenwart allezeit kund werden und erfahren lassen.

Blicket zurück auf den geringen, im Aeußerlichen so armseligen Anfang dieser Anstalt. Wer hat damals so herrlich durchgeholfen? Es war Jesus. Bei aller Armuth an zeitlichen Gütern war doch immer genug vorhanden. Wer sorgte nämlich dafür, daß das Wenige ausreichte? Es war Jesus. Bei aller Widerwärtigkeit ging das Werk doch herrlich voran. Wer sorgte nämlich dafür, daß alle Hindernisse glücklich überstanden wurden, ja, zum Segen dienten? Es war Jesus. Mochte den Professoren wohl manchmal ihre Arbeit sauer werden, Jesus gab ihnen dazu immer wieder neue Kraft. Mochte den Schülern wohl manchmal das Studiren schwer, ja, zuwider werden, Jesus gab ihnen dazu immer wieder neue Lust und Freudigkeit. Ja, mochte diese Anstalt vielfach von Menschen verachtet werden, so war und blieb sie doch eine hochbegnadigte Anstalt, eine Anstalt, in der Jesus selbst wohnte und wirkte und die er zu einer Stätte des Segens für viele machte.

Blicket ferner hin auf die mancherlei besonderen Nothstände, die im Laufe der Jahre über diese Anstalt kamen. Wer hat da immer wieder ausgeholfen? Es war Jesus. Wie oft wurden zum Beispiel treue und tüchtige Professoren der Anstalt genommen, sei es durch Wegberufung, oder durch den Tod, oder auf andere Weise! Wer hat dann aber immer wieder andere Männer willig und tüchtig gemacht, unsern Jünglingen als Lehrer hier zu dienen? Es war Jesus. Wie oft wurde der Raum zu klein, bald zum Studiren, bald zum Wohnen, bald zum Wirthschaften! Wer hat dann aber immer wieder die nöthigen Mittel gegeben, um mehr Raum zu schaffen? Es war Jesus. Auch wurde diese Anstalt von schweren Unglücksfällen heimgesucht. Denkt zum Beispiel daran, wie im Jahre 1869 das damalige Hauptgebäude in wenigen Stunden durch Feuer fast gänzlich zerstört wurde. Wer erweckte aber damals die Herzen unserer lieben Christen zum reichlichen Helfen und Geben, daß alle Noth in kurzer Zeit beseitigt und das Gebäude wieder hergestellt wurde? Es war Jesus.

Blicket ferner hin auf das Leben und Verhalten der Schüler in der Anstalt. Wo so viele junge Leute, wie hier, zusammenwohnen, da liegen auch der Teufel und das eigene böse Fleisch beständig auf der Lauer, um Schaden und Verderben anzurichten. Wer hat aber die vielen Jahre hindurch in Gnaden über unsere Schüler gewaltet, daß Gottes Wort die Zucht und Herrschaft unter ihnen behalten hat? Es war Jesus. Und wenn auch manchmal grobe Sünden und Ausschreitungen unter den Schülern vorkamen, was hat doch Jesus immer wieder gethan? Er hat Gnade zur Buße gegeben, oder er hat doch dafür gesorgt, daß unsere Anstalt keinen bleibenden Schaden von jenen Sünden und Ausschreitungen davongetragen hat.

Ja, meine Theuren, „der Herr ist noch und nimmer nicht von seinem Volk geschieden“. Das gilt auch von unserer theuren Concordia. Jesus war von Anfang an mit dieser Anstalt. Und er ist bis heute mit seiner Gnade und Hülfe, mit seinem Schutz und Segen bei ihr geblieben. Wie getrost können wir aber aus dieser gemachten Erfahrung den Schluß ziehen, daß er auch ferner bei ihr bleiben werde! Wie getrost und zuversichtlich können wir darum zu ihm beten: „Herr Jesu, bleibe bei uns!“

Doch, meine Lieben, was uns vor allen Dingen zu dieser Bitte: „Herr Jesu, bleibe bei uns!“ getrost und fröhlich macht, ist dieses, daß wir es aus seinem Wort so gewiß wissen, daß er diese unsere Bitte erfüllen will. Die Bitte der Emmausjünger: „Bleibe bei uns!“ wurde sogleich erfüllt. Denn es heißt in unserm Text: „Und er ging hinein, bei ihnen zu bleiben.“ Und obwohl Jesus ihnen bald darauf seine sichtbare Gegenwart entzog, so erschien er ihnen doch an demselben Abend wieder zu Jerusalem in sichtbarer Gestalt. Und wenige Wochen später gab er ihnen und allen seinen gläubigen Jüngern die ausdrückliche Verheißung: „Siehe, ich bin bei euch alle Tage, bis an der Welt Ende.“

O wie getrost und zuversichtlich können nun auch wir beten: „Herr Jesu, bleibe bei uns!“ Denn seine theure Verheißung: „Siehe, ich bin bei euch alle Tage, bis an der Welt Ende“ gilt auch uns. Ja, so gewiß Jesus es verheißen hat, so gewiß will und wird er mit uns und darum auch mit dieser unserer Anstalt sein und bleiben. Solange sein Wort hier im Schwange geht, so lange wird er selbst durch sein Wort hier zugegen sein.

Ja, Jesus will auch ferner das Haupt, der Herr und König unserer Anstalt bleiben. Er will, wie bisher, so auch in Zukunft für das geistliche und leibliche Wohl unserer Anstalt sorgen. Er will ihr Professoren und Lehrer geben; er will die zur Unterhaltung nöthigen Mittel darreichen; er will hier die Herzen lenken und regieren, daß sie ihm dienen. — Jesus will sich auch ferner in unserer Anstalt als Segenspender beweisen. Er will zur Arbeit der Professoren Frucht und Gedeihen geben. Und er will das Studiren der Schüler wohlgelingen lassen zu ihrem eigenen
\quellseite{4}
Segen und zum Segen vieler ihrer Mitmenschen. — Jesus will sich hier auch ferner als Berather und Helfer beweisen. Er will hier in dunklen Stunden den rechten Weg weisen. Er will hier die Muthlosen und Verzagten mit neuem Muth erfüllen, die Schwachen stärken, die müden Kniee aufrichten, den Irrenden zurechthelfen. — Jesus will sich hier auch ferner als Heiland beweisen. Er will hier mit seiner Gnade und seinem ganzen Verdienst zugegen sein, so daß Professoren und Schüler in ihrer Sündennoth zu ihm fliehen und Gnade und Vergebung bei ihm finden können. Allen soll hier durch Jesum der Weg zum Vaterherzen Gottes und in den Himmel offen stehen.

Sehet, das alles liegt in der Verheißung des Herrn Jesu: „Siehe, ich bin bei euch alle Tage, bis an der Welt Ende.“ Er selbst mit allem, was er ist und hat, will in unserer Anstalt wohnen, walten und wirken. Und darauf können und sollen wir uns so gewiß verlassen, so gewiß seine Verheißung wahr ist und nicht lügen kann. O wie getrost und zuversichtlich können wir darum heute den Herrn Jesum bitten, daß er auch ferner dieser Anstalt seine gnadenreiche Gegenwart erhalte. Wie getrost und zuversichtlich können wir heute bitten: „Herr Jesu, bleibe bei uns!“ —

Doch, meine Lieben, die Bitte: „Herr Jesu, bleibe bei uns!“ können wir nicht nur getrost und zuversichtlich thun, sondern zu dieser Bitte haben wir auch hohe Ursache. Und das wollen wir nun zweitens erwägen.

2.

Die Emmausjünger hatten unterwegs ihre eigene große Schwachheit und Gebrechlichkeit kennen gelernt. Sie hatten erfahren, was für unwissende, irrende, elende Menschen sie ohne Jesum waren. Sie hatten erfahren, daß sie ohne Jesum nichts vermochten. Und darum wollten sie Jesum nicht von sich scheiden sehen. Darum baten sie: „Bleibe bei uns!“

„Herr Jesu, bleibe bei uns!“ zu dieser Bitte haben auch wir hohe Ursache, wenn wir erstlich auf uns selbst sehen. Denn wie steht es mit uns, soviel auf uns selbst ankommt? Das sagt uns unser Heiland mit den Worten: „Ohne mich könnt ihr nichts thun.“ Wir müssen mit Luther bekennen:

\begin{verse}
Mit unsrer Macht ist nichts gethan,\\
Wir sind gar bald verloren.
\end{verse}

Ach ja, ohne Jesum sind wir ganz und gar hülflos, elend und verloren.

Was wäre darum auch unsere Lehranstalt ohne Jesum? Sie wäre ein trauriges Menschenwerk ohne allen Werth in Gottes Augen. Ohne Jesum fehlte dieser Anstalt das rechte Haupt, das Gott selbst über sein Volk gesetzt hat. Ohne Jesum fehlte dieser Anstalt der rechte Grund- und Eckstein, worauf Gott selbst sein Reich gebaut hat. Ohne Jesum würden hier keine Christen erzogen, sondern Heiden, keine Seelenhirten, sondern Seelenverderber. O gewiß, wir haben hohe Ursache, von Herzen zu bitten: „Herr Jesu, bleibe bei uns!“ Denn nur, wenn Jesus in und mit dieser Anstalt bleibt, kann Gott sein Werk hier haben und sein Reich hier bauen.

Wie stände es ferner mit den Professoren dieser Anstalt ohne Jesum? Alle ihre noch so fleißige Arbeit wäre ohne Jesum vergebens. Ohne Jesum könnten sie ihre Schüler nicht recht lehren und nicht recht erziehen. Ja, ohne Jesum würden sie mit ihren Schülern falsche Bahnen betreten und verderbliche Irrwege einschlagen. Alles, alles darum, was ihnen zu ihrem Amt nöthig ist, muß Jesus geben. Nur von Jesu kommt ihnen die rechte Weisheit zum Lehren und Erziehen; nur von Jesu Segen zu ihrer Arbeit; nur von Jesu Geduld mit den Schwachen; nur von Jesu Muth und Freudigkeit, wenn das Herz verzagt werden will; nur von Jesu Trost und Gnade, wenn die eigenen Sünden und Gebrechen sie anfechten. Ach ja, unsere lieben Professoren haben den Herrn Jesum sehr nöthig. Welch hohe Ursache haben wir darum zu bitten: „Herr Jesu, bleibe bei uns!“

Wie stände es ferner mit den Schülern dieser Anstalt ohne Jesum? Auch ihre Arbeit, ihr Lernen und Studiren, wäre ohne Jesum vergebens. Ja, ohne Jesum fehlte ihnen schon die rechte Lust und Liebe zur Arbeit. Ohne Jesum würden sie ihre Studien auch nicht in der rechten Gesinnung treiben; nicht aus Liebe zum Heiland, sondern um Menschen zu gefallen; nicht, um ihrem Nächsten zu dienen, sondern um ihre eigene Ehre zu suchen. Auch den Schülern muß darum alles von Jesu kommen, was sie nöthig haben. Nur von Jesu kommt ihnen Lust, Kraft und Ausdauer zum Studiren. Nur von Jesu kommt ihnen Segen und Erfolg zu ihrer Arbeit. Nur von Jesu lernen sie treuen Gehorsam, rechte Demuth, wahre Gottseligkeit, so daß von ihnen gesagt werden kann: Sie nehmen zu an Weisheit und Gnade bei Gott und den Menschen. Ach ja, auch unsere Schüler haben den Herrn Jesum sehr nöthig. Welch hohe Ursache haben wir darum zu bitten: „Herr Jesu, bleibe bei uns!“

Doch, meine Lieben, zu dieser Bitte haben wir hohe Ursache, nicht nur wenn wir auf uns selbst sehen, sondern auch wenn wir auf unsere Feinde sehen. Die Emmausjünger waren jetzt glückliche und selige Leute, da sie Jesum bei sich hatten. Sie wußten aber auch, daß sie viele Feinde hatten, vor denen sie keinen Augenblick sicher waren. So hatten sie denn auch deswegen hohe Ursache, Jesum zu bitten: „Bleibe bei uns!“ Auch bei uns steht es so. Besonders jetzt, da es Abend werden will, da das Ende der Welt naht, sind die Christen allenthalben von Feinden umgeben, die uns zu schaden und zu verderben trachten. Und das gilt, wie von uns, so auch von unserer Lehranstalt.

Welch einen grimmigen Feind hat unsere Anstalt am Teufel, weil dieser weiß, daß durch sie seinem Reich großer Abbruch geschieht. Er trachtet hier bald auf die eine, bald auf die andere Weise, Schaden anzurichten. Er versucht und reizt hier zum Irrthum, Unglauben und Abfall. Er versucht und reizt hier zu Sünden und zum gottlosen Leben. Und das alles gerade jetzt, da er weiß, daß er nur noch wenig Zeit hat. Und der Teufel ist ein gar schlimmer Feind.

\begin{verse}
Groß Macht und viel List\\
Sein grausam Rüstung ist;\\
Auf Erd ist nicht seins Gleichen.
\end{verse}

Nun sagt, wer kann hier helfen? Wer kann unsere Anstalt beschützen, daß der Teufel hier nicht Schaden und Verderben anrichte? Das kann Jesus allein.

Welch eine giftige Feindin unserer Anstalt ist ferner die Welt. Sie haßt diese Anstalt, weil es eine christliche Anstalt ist. Sie verspottet die Schüler, weil sie nicht große Männer dieser Welt, sondern Prediger des Evangeliums werden wollen. Sie sucht hier mit ihrem Weltwesen einzudringen und dadurch Unheil anzurichten. Wer kann nun auch hier helfen und der Feindschaft der Welt wehren? Das kann Jesus allein.

\quellseite{5}
Welch ein mächtiger Feind ist endlich das eigene böse Fleisch. Dieses will fort und fort unsere Gemeinden träge machen und ihnen das Interesse für unsere Anstalt nehmen. Dieses setzt auch unsern Professoren auf mancherlei Weise zu und ist gar mächtig in den Schülern, sie vom Guten abzuhalten und zum Bösen anzutreiben. Wer kann nun auch hier helfen und dem bösen Fleisch die Herrschaft nehmen? Das kann Jesus allein.

O wahrlich, meine Lieben, so haben wir denn, auch wenn wir auf die Feinde unserer Anstalt sehen, hohe Ursache, zu beten und zu bitten: „Herr Jesu, bleibe bei uns!“

Wohlan, meine theuren Mitchristen, da wir wissen, daß unsere Concordia nur mit Jesu wohl daran ist, daß Jesus hier alles thun und alles vollbringen muß von Anfang bis zu Ende, wie, wollen und sollen wir ihn da nicht herzlich bitten, daß er, wie bisher, so auch ferner bei und mit ihr bleibe? O darum Herzen und Augen nach oben! Und mögen auch unsere Lippen nicht müde werden zu bitten und zu beten: „Herr Jesu, bleibe bei uns!“ Er will es thun — er wird es thun. Amen!
